\subsectionaddtolist{الف}

تعریف بخش زوج سیگنال $h(t)$ به صورت زیر است:
\[ h_e(t) = \frac{h(t) + h(-t)}{2} \]

با توجه به این که سیستم یک سیستم LTI {علی} است، طبق تعریف سیستم‌های علّی، پاسخ ضربه سیستم برای زمان‌های منفی باید صفر باشد. یعنی:
\[ h(t) = 0 \quad t < 0 \]

حال این شرط را در دو حالت بررسی می‌کنیم:

\begin{enumerate}
    \item {برای حالت $t > 0$:}
    چون $t > 0$ است، پس $-t < 0$ خواهد بود. در نتیجه با توجه به خاصیت علّیت داریم $h(-t) = 0$. با جایگذاری این مقدار در رابطه بخش زوج:
    \[ h_e(t) = \frac{h(t) + 0}{2} = \frac{h(t)}{2} \implies h(t) = 2h_e(t) \quad (t > 0) \]
    
    \item {برای حالت $t < 0$:}
    همانطور که گفته شد، به دلیل علّیت سیستم، مقدار پاسخ ضربه برای زمان‌های منفی صفر است:
    \[ h(t) = 0 \quad (t < 0) \]
\end{enumerate}

بنابراین، کل سیگنال $h(t)$ را می‌توان برحسب $h_e(t)$ به صورت زیر بازنویسی کرد (با فرض عدم وجود تکینگی در $t=0$):
\[ h(t) = 2h_e(t)u(t) \]
که در آن $u(t)$ تابع پله واحد است.

\subsectionaddtolist{ب}
فرض شده است که:
\[ \mathcal{F}\{h_e(t)\} = \Re\{H(j\omega)\} = H_R(j\omega) = \begin{cases} 1, & |\omega| < W \\ 0, & |\omega| > W \end{cases} \]

ابتدا با استفاده از تبدیل فوریه معکوس، $h_e(t)$ را به دست می‌آوریم:
\[ h_e(t) = \frac{1}{2\pi} \int_{-\infty}^{\infty} H_R(j\omega) e^{j\omega t} d\omega = \frac{1}{2\pi} \int_{-W}^{W} 1 \cdot e^{j\omega t} d\omega \]

اگر $t \neq 0$ باشد:
\[ h_e(t) = \frac{1}{2\pi} \left[ \frac{e^{j\omega t}}{jt} \right]_{-W}^{W} = \frac{1}{2\pi j t} \left( e^{jWt} - e^{-jWt} \right) \]
با استفاده از رابطه اویلر $\sin(Wt) = \frac{e^{jWt} - e^{-jWt}}{2j}$ داریم:
\[ h_e(t) = \frac{\sin(Wt)}{\pi t} \]

اکنون با استفاده از نتیجه‌ی بخش (الف)، پاسخ ضربه $h(t)$ برابر است با:
\[ h(t) = 2h_e(t)u(t) = \frac{2\sin(Wt)}{\pi t} u(t) \]

\subsectionaddtolist{ج}
 بله، این سیستم علّی است.
 طبق رابطه به دست آمده در بخش (ب)، پاسخ ضربه سیستم یعنی $h(t) = \frac{2\sin(Wt)}{\pi t} u(t)$ به دلیل ضرب شدن در تابع پله واحد $u(t)$، برای تمام زمان‌های منفی دقیقاً برابر با صفر است. از آنجا که شرط اصلی علّیت یک سیستم LTI مستقل بودن خروجی از مقادیر آینده ورودی (یا به عبارتی صفر بودن پاسخ ضربه در زمان‌های منفی) است، این سیستم کاملاً علّی می‌باشد.
